% GENERATED FILE. Edit canonical lessons, then run npm run generate:books.
% canonical-source-hash: fnv1a64:7b8676f28a3c5c81
% canonical-lessons: LA-C300-tabellarius, LA-C300-praeco, LA-C300-mandatum, LA-C300-legatio, LA-C300-renuntio

\chapter{Tabellārius et praecō}
\label{ch:la-tabellarius-et-praeco}
\hlchaptermodality{300}

\begin{chapteropening}
\textbf{By the end of this chapter:} \emph{I can talk about how news travelled: the letter-carrier, the herald, the instructions, the embassy and bringing word back.}

Complete the last lesson of chapter 300: I can talk about how news travelled: the letter-carrier, the herald, the instructions, the embassy and bringing word back.
\end{chapteropening}

\section[tabellārius, tabellāriī]{tabellārius, tabellāriī — a letter-carrier, a courier}

\label{lesson:LA-C300-tabellarius}



\begin{quote}
Before the new one: say the Latin for to report, then the Latin for I say.
\end{quote}

\subsection*{You'll want to know: tabellārius, tabellāriī}
\textbf{tabellārius, tabellāriī} — \textquotedblleft{}a letter-carrier, a courier\textquotedblright{}.

\begin{quote}
Tabellārius epistulam patris ad nōs portāvit.
\end{quote}

\begin{quote}
\emph{The courier carried father's letter to us.}
\end{quote}

Masculine. Romans with money sent their own slaves as couriers; letters went on wax \emph{tabellae}.

\begin{cousinweb}
\textbf{tabellārius} is built on \emph{tabella}, \textquotedblleft{}a small writing tablet, a letter\textquotedblright{} (Lewis \& Short): the man who carries the tablets.
\end{cousinweb}

\begin{grammarlens}[title={What you've built}]
One more word for telling, explaining and describing at B1.
\end{grammarlens}

\subsection*{Your turn}
\begin{itemize}
\raggedright
  \item \emph{Say it:} \emph{tabellārius, tabellāriī}
  \item \emph{Say it:} \emph{tabellārius, tabellāriī}, once more
  \item \emph{Say it:} a sentence of your own with \emph{tabellārius, tabellāriī}
  \item \emph{Recall:} say the Latin for to report, then the Latin for I say, then say \emph{tabellārius, tabellāriī} again
\end{itemize}

\begin{tcolorbox}[breakable,colback=teal!4,colframe=teal!35!black,title={Before you move on}]
What does tabellārius, tabellāriī mean? (\textquotedblleft{}A letter-carrier, a courier\textquotedblright{}.) Say it once more.
\end{tcolorbox}
\section[praecō, praecōnis]{praecō, praecōnis — a herald, a crier}

\label{lesson:LA-C300-praeco}



\begin{quote}
Before the new one: say the Latin for I say, then the Latin for a letter-carrier.
\end{quote}

\subsection*{You'll want to know: praecō, praecōnis}
\textbf{praecō, praecōnis} — \textquotedblleft{}a herald, a crier\textquotedblright{}.

\begin{quote}
Praecō in forō victōriam nūntiāvit.
\end{quote}

\begin{quote}
\emph{The herald announced the victory in the forum.}
\end{quote}

Masculine, 3rd declension. He called for silence at trials, announced the lots at auctions and proclaimed winners at the games.

\begin{cousinweb}
\textbf{praecō} is usually traced to \emph{prae}, \textquotedblleft{}before, in front\textquotedblright{}, and \emph{vocō}, \textquotedblleft{}to call\textquotedblright{}: one who calls out in front. Lewis \& Short give the word without a settled origin.
\end{cousinweb}

\begin{grammarlens}[title={What you've built}]
One more word for telling, explaining and describing at B1.
\end{grammarlens}

\subsection*{Your turn}
\begin{itemize}
\raggedright
  \item \emph{Say it:} \emph{praecō, praecōnis}
  \item \emph{Say it:} \emph{praecō, praecōnis}, once more
  \item \emph{Say it:} a sentence of your own with \emph{praecō, praecōnis}
  \item \emph{Recall:} say the Latin for I say, then the Latin for a letter-carrier, then say \emph{praecō, praecōnis} again
\end{itemize}

\begin{tcolorbox}[breakable,colback=teal!4,colframe=teal!35!black,title={Before you move on}]
What does praecō, praecōnis mean? (\textquotedblleft{}A herald, a crier\textquotedblright{}.) Say it once more.
\end{tcolorbox}
\section[mandātum, mandātī]{mandātum, mandātī — an instruction, an order}

\label{lesson:LA-C300-mandatum}



\begin{quote}
Before the new one: say the Latin for a letter-carrier, then the Latin for a herald.
\end{quote}

\subsection*{You'll want to know: mandātum, mandātī}
\textbf{mandātum, mandātī} — \textquotedblleft{}an instruction, an order\textquotedblright{}.

\begin{quote}
Lēgātus mandāta imperātōris memoriā tenēbat.
\end{quote}

\begin{quote}
\emph{The envoy kept the general's instructions in his memory.}
\end{quote}

Neuter, from \emph{mandō} (taught), 'to entrust'. The plural \emph{mandāta} is the message an envoy carried. English \emph{mandate}.

\begin{cousinweb}
\textbf{mandātum} is the neuter participle of \emph{mandō}, \textquotedblleft{}to entrust, to order\textquotedblright{} (Lewis \& Short): what has been entrusted. English \emph{mandate} comes from it.
\end{cousinweb}

\begin{grammarlens}[title={What you've built}]
One more word for telling, explaining and describing at B1.
\end{grammarlens}

\subsection*{Your turn}
\begin{itemize}
\raggedright
  \item \emph{Say it:} \emph{mandātum, mandātī}
  \item \emph{Say it:} \emph{mandātum, mandātī}, once more
  \item \emph{Say it:} a sentence of your own with \emph{mandātum, mandātī}
  \item \emph{Recall:} say the Latin for a letter-carrier, then the Latin for a herald, then say \emph{mandātum, mandātī} again
\end{itemize}

\begin{tcolorbox}[breakable,colback=teal!4,colframe=teal!35!black,title={Before you move on}]
What does mandātum, mandātī mean? (\textquotedblleft{}An instruction, an order\textquotedblright{}.) Say it once more.
\end{tcolorbox}
\section[lēgātiō, lēgātiōnis]{lēgātiō, lēgātiōnis — an embassy}

\label{lesson:LA-C300-legatio}



\begin{quote}
Before the new one: say the Latin for a herald, then the Latin for an instruction.
\end{quote}

\subsection*{You'll want to know: lēgātiō, lēgātiōnis}
\textbf{lēgātiō, lēgātiōnis} — \textquotedblleft{}an embassy\textquotedblright{}.

\begin{quote}
Lēgātiō ab hostibus vēnit et pācem petīvit.
\end{quote}

\begin{quote}
\emph{An embassy came from the enemy and asked for peace.}
\end{quote}

Feminine, 3rd declension. The \emph{lēgātus} (taught) led it. English \emph{legation}.

\begin{cousinweb}
\textbf{lēgātiō} is built on \emph{lēgō}, \textquotedblleft{}to send with a commission\textquotedblright{} (EtymDB): the sending of envoys, and the envoys themselves. English \emph{legation} comes from it.
\end{cousinweb}

\begin{grammarlens}[title={What you've built}]
One more word for telling, explaining and describing at B1.
\end{grammarlens}

\subsection*{Your turn}
\begin{itemize}
\raggedright
  \item \emph{Say it:} \emph{lēgātiō, lēgātiōnis}
  \item \emph{Say it:} \emph{lēgātiō, lēgātiōnis}, once more
  \item \emph{Say it:} a sentence of your own with \emph{lēgātiō, lēgātiōnis}
  \item \emph{Recall:} say the Latin for a herald, then the Latin for an instruction, then say \emph{lēgātiō, lēgātiōnis} again
\end{itemize}

\begin{tcolorbox}[breakable,colback=teal!4,colframe=teal!35!black,title={Before you move on}]
What does lēgātiō, lēgātiōnis mean? (\textquotedblleft{}An embassy\textquotedblright{}.) Say it once more.
\end{tcolorbox}
\section[renūntiō, renūntiāre]{renūntiō, renūntiāre — to bring back word, to report}

\label{lesson:LA-C300-renuntio}



\begin{quote}
Before the new one: say the Latin for an instruction, then the Latin for an embassy.
\end{quote}

\subsection*{You'll want to know: renūntiō, renūntiāre}
\textbf{renūntiō, renūntiāre} — \textquotedblleft{}to bring back word, to report\textquotedblright{}.

\begin{quote}
Servus ad dominum rediit et omnia renūntiāvit.
\end{quote}

\begin{quote}
\emph{The slave went back to his master and reported everything.}
\end{quote}

\emph{renūntiō, renūntiāre, renūntiāvī, renūntiātum}: \emph{re-} 'back' + \emph{nūntiō} (taught). Also 'to declare someone elected', and 'to give up, to renounce'.

\begin{cousinweb}
\textbf{renūntiō} is built on \emph{nūntiō}, \textquotedblleft{}to announce\textquotedblright{}, with \emph{re-}, \textquotedblleft{}back\textquotedblright{} (Lewis \& Short): to carry word back. English \emph{renounce} comes from another sense, \textquotedblleft{}to call off, to give up\textquotedblright{}.
\end{cousinweb}

\begin{grammarlens}[title={What you've built}]
One more word for telling, explaining and describing at B1.
\end{grammarlens}

\subsection*{Your turn}
\begin{itemize}
\raggedright
  \item \emph{Say it:} \emph{renūntiō, renūntiāre}
  \item \emph{Say it:} \emph{renūntiō, renūntiāre}, once more
  \item \emph{Say it:} a sentence of your own with \emph{renūntiō, renūntiāre}
  \item \emph{Recall:} say the Latin for an instruction, then the Latin for an embassy, then say \emph{renūntiō, renūntiāre} again
\end{itemize}

\begin{tcolorbox}[breakable,colback=teal!4,colframe=teal!35!black,title={Before you move on}]
What does renūntiō, renūntiāre mean? (\textquotedblleft{}To bring back word, to report\textquotedblright{}.) Say it once more.
\end{tcolorbox}
